% 실험 방법 섹션 (main.tex에 \input으로 포함)
% 이 파일은 \chapter{실험} 아래 \section{실험 설정}을 대체함

\section{실험 설정}

본 절에서는 제안하는 SE-gated cascade 방법의 성능 평가에 사용한 실험 설정을 설명한다.

\subsection{데이터셋}

TruthfulQA 데이터셋의 검증 분할(validation split)에서 무작위로 선택한 200개 샘플을 사용하였다. 이 데이터셋은 Lin 등(2022)이 인간이 흔히 가지는 오개념(misconception)을 유도하는 질문들로 구성한 것으로, LLM의 환각 탐지 연구에 널리 사용된다.

\begin{table}[htbp]
\centering
\caption{실험 데이터셋 통계}
\label{tab:dataset}
\begin{tabular}{@{}lr@{}}
\toprule
항목 & 값 \\
\midrule
실험 샘플 수 & 200 \\
환각 샘플 & 164 (82.0\%) \\
정상 샘플 & 36 (18.0\%) \\
\bottomrule
\end{tabular}
\end{table}

각 질문에는 복수의 정답(correct\_answers)과 오답(incorrect\_answers) 레이블이 부여되어 있으며, LLM 응답이 정답 중 하나와 부분 일치하면 정상(0), 아니면 환각(1)으로 분류하였다.


\subsection{실험 파이프라인}

각 질문 $q$에 대해 Qwen2.5-3B-Instruct로 $K=5$개의 응답을 샘플링하였다. 프롬프트는 ``Answer the question concisely in one sentence.''와 질문 본문으로 구성했으며, 생성 파라미터는 temperature 0.7, max tokens 50, do sample True로 설정하였다. 각 응답에 대해 생성 토큰의 원시 logit 값을 저장하여 Semantic Energy 계산에 사용하였다.

샘플링된 응답들은 DeBERTa-large-mnli를 사용해 의미 클러스터로 묶었다. 두 응답이 양방향 entailment 관계에 있으면 의미적으로 동등하다고 보고 같은 클러스터에 배치하였다.

\subsubsection{지표 계산}

\paragraph{Semantic Entropy (SE)}
클러스터 분포로부터 Shannon entropy를 계산한다:
\begin{equation}
    SE = -\sum_{c \in C} p(c) \log p(c), \quad p(c) = \frac{|c|}{K}
\end{equation}

\paragraph{Semantic Energy}
각 응답의 원시 logit 값으로부터 에너지를 계산한다:
\begin{equation}
    Energy = -\frac{1}{K} \sum_{i=1}^{K} \frac{1}{T_i} \sum_{t=1}^{T_i} \ell_{i,t}
\end{equation}
여기서 $\ell_{i,t}$는 $i$번째 응답의 $t$번째 토큰 logit이고, $T_i$는 응답 길이이다.

SE-gated cascade 평가는 위 두 점수를 모두 계산한 뒤 질문별로 하나의 최종 탐지 점수를 선택한다. 반복 샘플링 결과가 여러 의미 클러스터로 나뉘면($|C|\geq 2$) SE를 사용하고, 모든 응답이 하나의 클러스터에 모이면($|C|=1$) Energy를 사용한다. 따라서 cascade 점수는 별도의 새 지표가 아니라, SE가 정보를 갖는 구간과 그렇지 않은 구간을 나누어 기존 두 신호를 선택적으로 사용하는 값이다.


\subsubsection{환각 판정 및 평가}

$K$개 응답 중 하나라도 정답(correct\_answers)과 부분 일치하면 정상(0), 모든 응답이 오답이면 환각(1)으로 레이블링하였다. 각 질문은 SE, Energy, 또는 SE-gated cascade 점수를 받으며, 이 점수가 높은 사례일수록 환각 가능성이 큰 것으로 정렬된다. 탐지 성능은 AUROC와 AUPRC로 측정하였다.


\subsection{모델 설정}

\begin{table}[htbp]
\centering
\caption{실험에 사용된 모델}
\label{tab:models}
\begin{tabular}{@{}lll@{}}
\toprule
역할 & 모델 & 비고 \\
\midrule
LLM (응답 생성) & Qwen2.5-3B-Instruct & HuggingFace Transformers \\
NLI (클러스터링) & DeBERTa-large-mnli & 3-way classification \\
\bottomrule
\end{tabular}
\end{table}

\paragraph{LLM 선택 근거}
Qwen2.5-3B-Instruct는 명령어 튜닝을 거친 모델로, 질의응답에 적합하며 3B 파라미터로 실험 효율성을 확보하였다. 더 크거나 다양한 LLM 설정에서의 비교는 향후 연구에서 다룬다.

\paragraph{NLI 모델 선택 근거}
DeBERTa-large-mnli는 MNLI 벤치마크에서 90\% 이상의 정확도를 달성한 모델로, Semantic Entropy 원 논문(Farquhar et al., 2024)에서도 사용되었다.


\subsection{실험 환경}

\begin{table}[htbp]
\centering
\caption{실험 환경}
\label{tab:environment}
\begin{tabular}{@{}ll@{}}
\toprule
항목 & 설정 \\
\midrule
GPU & NVIDIA RTX 5090 (32GB) \\
CUDA & 12.1 \\
Python & 3.11 \\
PyTorch & 2.1.0 \\
Transformers & 4.36.0 \\
\bottomrule
\end{tabular}
\end{table}


\subsection{재현성}

본 아카이브에는 논문 원문, 빌드 가능한 LaTeX 소스, 그리고 논문에 사용한 결과 JSON이 포함된다. 다만 실험 실행 코드와 실행용 패키지 소스 코드는 포함하지 않으므로, 본 저장소만으로 원 실험을 재실행하기에는 제한이 있다.
결과 검증은 공개된 JSON의 설정값과 집계값을 기준으로 수행하였다.
